% !Mode:: "TeX:UTF-8"
\begin{conclusions}

本文围绕PyTorch-FSDP分布式训练框架下的通信优化问题，系统研究了通信数据压缩技术的实现方法、优化策略和性能影响。主要工作与研究成果总结如下：

（1）通过构建单卡与多卡GPT-2模型对比实验，利用PyTorch Profiler对FSDP训练过程进行了细粒度性能分析，定量识别了通信瓶颈。实验表明在GPT2-large模型双卡训练中，通信开销占单步耗时的61.2\%，且Reduce-Scatter操作位于反向传播关键路径上，阻塞效应显著，从而确立了后向梯度通信作为压缩优化首要目标的必要性。

（2）设计了基于FSDP通信钩子机制的模块化通信压缩框架，将压缩策略与训练逻辑解耦，实现了量化类（INT8、QSGD、Natural Compression、1-bit Seide）和稀疏化类（Top-k、Random-k、Threshold-v、Sketched-SGD）共9种压缩算法的统一集成，支持算法的热插拔与组合使用。

（3）在统一实验配置下对各算法进行了系统性对比评估。实验结果表明：INT8量化实现约4倍压缩比，训练吞吐提升约35\%，Loss仅上升3.8\%；Top-k稀疏化（1\%稀疏率）在保持良好收敛质量的同时达到最高吞吐；梯度相对L2误差控制在1.0以内时训练收敛质量可接受。INT8量化、Top-k和Threshold-v稀疏化在压缩比、通信时间和收敛质量的综合权衡上表现最为均衡。

本文的研究仍存在以下不足，可作为未来工作的方向：当前实验环境局限于单机双卡，通信瓶颈尚未充分暴露，后续应在多机多卡环境下进一步验证各压缩算法的可扩展性；可以探索将量化与稀疏化方法进行组合（如先Top-k再INT8量化），以实现更高的综合压缩比；此外，自适应压缩策略（根据训练阶段动态调整压缩强度）也是值得深入研究的方向。

\end{conclusions}
